
\documentclass[10pt,a4paper]{article}
\usepackage[utf8]{inputenc}
\usepackage{booktabs}
\usepackage[a4paper,left=3cm,right=2cm]{geometry}
\usepackage{amsmath}
\usepackage[thmmarks, amsmath, thref]{ntheorem}
\usepackage{amsfonts}
\usepackage{amssymb}
\usepackage{fancyhdr}
\pagestyle{fancy}
\usepackage{anysize}
\usepackage{graphicx}
\usepackage{float}
\usepackage{hyperref}
\usepackage{enumitem}
\usepackage{verbatim}
\newtheorem{theorem}{Theorem}
\theoremsymbol{\ensuremath{\blacksquare}}
\newtheorem*{solution}{Solución:}

\usepackage{xcolor}
\usepackage{listings}
% Margenes y formalidades

\textwidth 6.25in % Margen Derecho y ancho del texto
\textheight 8.5in
\oddsidemargin 0in % Margen Izquierdo
\headheight 0in % Largo del texto
\leftmargin 1in
\parindent 0em % Salto entre la indentación (Sangría)https://www.overleaf.com/4519127669sdpptkhvsbfn
\topmargin -0.5in % Margen Superior
\parskip 0.5ex % Salto entre parrafos
\baselineskip 1pt



\lhead{
\includegraphics[scale=0.08]{Figuras/logoUSM.png}} % texto izquierda de la cabecera
%\chead{}                                      % texto centro de la cabecera
\rhead{Universidad Técnica Federico Santa María\\Ingeniería Comercial } % número de página a la derecha
\renewcommand{\headrulewidth}{0.4pt} % grosor de la línea de la cabecera
\renewcommand{\footrulewidth}{0.4pt} % grosor de la línea del pie

\begin{document}
\vspace{5cm}



\begin{center}
\begin{large}
\textbf{\\CERTAMEN 3:}\textbf{ ECONOMETRÍA ICS-294}
\end{large}\\
\vspace{0.2cm}
II Semestre 2025\\
Profesores: Jocelyn Tapia Stefanoni, Franco Antonucci Fernández y Ricardo Pruneda Urz\'ua\\
%Ayudante: 
\end{center}

\textbf{Instrucciones:} 

\begin{itemize}

    \item Usted dispone de 90 minutos  para intentar los 100 puntos del certamen.
    \item Cuide la presentación y redacción de sus preguntas.
    \item No puede utilizar celular, computador o tablet, pero sí una calculadora TI.
     
\end{itemize}

\section{RdeA9.1 Comentes y preguntas cortas (40 puntos)}
{\bf Elija y responda solo 4 de las siguientes preguntas.}
\begin{enumerate}
    \item 
Considere la base de datos \texttt{bwght} de Wooldridge, que contiene información sobre recién nacidos
y características de sus madres durante el embarazo. Se define la variable dependiente
$lowbirth_i$ como una variable dicotómica que toma valor 1 si el bebé $i$ nace con bajo peso
(por ejemplo, menos de 2500 gramos), y 0 en caso contrario.

A partir de la base, se ha estimado el siguiente modelo de probabilidad lineal (MPL):

\[
\widehat{lowbirth}_i
= 0.20 + 0.02\, cigs_i - 0.004\, faminc_i,
\]

donde:
\begin{itemize}
    \item $cigs_i$: número de cigarrillos que la madre fuma por día durante el embarazo,
    \item $faminc_i$: ingreso familiar medido en miles de dólares anuales.
\end{itemize}

Responda:

\begin{enumerate}
    \item (2 puntos) Interprete el coeficiente $\hat\beta_0 = 0.20$.

\color{blue}
La interpretación del intercepto $\hat\beta_0 = 0.20$ es la siguiente:
    cuando $cigs_i = 0$ y $faminc_i = 0$, la probabilidad estimada de que el bebé
    nazca con bajo peso es
    \[
        \widehat{lowbirth} = 0.20,
    \]
    es decir, un 20\%. Representa la probabilidad base del modelo cuando las
    variables explicativas toman el valor cero.
\color{black}

    
    \item (2 puntos) Interprete el coeficiente asociado a la variable $cigs_i$.

    \color{blue}
El coeficiente asociado a $cigs_i$ es $0.02$. Esto significa que,
    manteniendo constante el ingreso familiar,
    \[
        \frac{\partial \widehat{lowbirth}}{\partial cigs} = 0.02,
    \]
    por lo que cada cigarrillo adicional que la madre fuma al día aumenta la
    probabilidad de bajo peso en 0.02, es decir, en 2 puntos porcentuales.
\color{black}
    \item (2 puntos) Interprete el coeficiente asociado a la variable $faminc_i$.

    \color{blue}
El coeficiente asociado a $faminc_i$ es $-0.004$. Esto implica que,
    manteniendo constante el consumo de cigarrillos,
    \[
        \frac{\partial \widehat{lowbirth}}{\partial faminc} = -0.004,
    \]
    por lo que un aumento de una unidad en el ingreso familiar (mil dólares)
    reduce la probabilidad de bajo peso en 0.004, es decir, en 0.4 puntos
    porcentuales.
\color{black}
    \item (2 puntos)  Calcule la probabilidad de bajo peso para un bebé cuya madre fuma 10 cigarrillos al día y tiene ingreso familiar de 20 (miles de dólares).

    \color{blue}
Sustituyendo $cigs_i = 10$ y $faminc_i = 20$ en la ecuación estimada:
    \[
        \widehat{lowbirth}
        = 0.20 + 0.02(10) - 0.004(20).
    \]
    Calculando:
    \[
        0.02(10) = 0.20, \qquad 0.004(20) = 0.08,
    \]
    \[
        \widehat{lowbirth} = 0.20 + 0.20 - 0.08 = 0.32.
    \]
    Por lo tanto, la probabilidad estimada es
    \[
        \widehat{lowbirth} = 0.32 \quad (32\%).
    \]

\color{black}
    \item (2 puntos)  ¿Qué problema del modelo de probabilidad lineal se evidencia al considerar valores extremos de las variables?

\color{blue}
El modelo de probabilidad lineal presenta el problema de que las
    predicciones no están restringidas al intervalo $[0,1]$. Para valores
    extremos de las variables, como un número muy alto de cigarrillos o un
    ingreso familiar muy bajo, el modelo puede entregar probabilidades
    negativas o superiores a 1, lo cual es imposible para una probabilidad.
\color{black}
\end{enumerate}
    
       \item Luego de estimar un modelo de regresión, se aplica el test de Breusch--Pagan para evaluar la presencia de heterocedasticidad en los residuos. El test arroja un \textit{p-valor} de 0.004.

\begin{enumerate}
  \item[a)] (3 puntos) ¿Qué implica este resultado respecto al cumplimiento del supuesto de homocedasticidad en el modelo estimado?
  \item[b)] (4 puntos) ¿Qué efectos tiene la heterocedasticidad sobre los errores estándar y sobre la validez de la inferencia estadística basada en el modelo?
  \item[c)] (3 puntos) ¿Qué métodos pueden utilizarse para corregir o mitigar este problema?
\end{enumerate}



{\color{blue}  \textbf{a) Interpretación del p-valor del test de Breusch--Pagan}

El test de Breusch--Pagan contrasta la hipótesis nula:

\[
H_0: \operatorname{Var}(\varepsilon_i)=\sigma^2 \quad \text{(homocedasticidad)}
\]

contra la alternativa:

\[
H_1: \operatorname{Var}(\varepsilon_i)=\sigma_i^2  \quad \text{(heterocedasticidad)}
\]

donde $z_i$ son variables explicativas o funciones de ellas.

En la práctica, el test se implementa estimando la siguiente \textit{ecuación auxiliar}:

\[
\hat{u}_i^2 = \gamma_0 + \gamma_1 x_{1i} + \gamma_2 x_{2i} + \cdots + \gamma_k x_{ki} + v_i ,
\]

donde $\hat{u}_i$ son los residuos del modelo original. Bajo $H_0$, se tiene que:

\[
\gamma_1 = \gamma_2 = \cdots = \gamma_k = 0.
\]

El estadístico del test se construye mediante:

\[
LM = n R^2_{\text{aux}},
\]

que bajo $H_0$ se distribuye asintóticamente como $\chi^2_k$.

(Con la siguiente parte  es suficiente para obtener todo el puntaje)

Dado que el p-valor obtenido es $0.004$, menor a niveles de significancia comunes (1\%, 5\%), se rechaza $H_0$. Por lo tanto, existe evidencia estadística de heterocedasticidad y el supuesto de varianza constante no se cumple en el modelo.


\textbf{b) Efectos de la heterocedasticidad}

La heterocedasticidad no afecta la insesgadez de los estimadores OLS, pero sí altera sus propiedades de eficiencia y la validez de la inferencia:

\begin{itemize}
    \item Los errores estándar OLS se vuelven \textit{inconsistentes}, pudiendo estar subestimados o sobrestimados.
    \item Como consecuencia, los estadísticos $t$ y $F$ son incorrectos.
    \item Los intervalos de confianza pueden ser demasiado estrechos o demasiado amplios.
    \item Se pueden cometer errores en pruebas de significancia: falsos positivos o falsos negativos.
\end{itemize}

Por ello, la heterocedasticidad compromete la validez de la inferencia estadística basada en el modelo.

\vspace{0.6cm}

\textbf{c) Métodos de corrección o mitigación}

Existen varias estrategias habituales para tratar la heterocedasticidad:

Basta con la primera para tener todo el puntaje.
\begin{enumerate}
    \item \textbf{Errores estándar robustos (White)}: permiten realizar inferencia válida sin cambiar los coeficientes OLS.
    \item \textbf{Mínimos Cuadrados Generalizados (GLS)}: más eficientes. 
     \end{enumerate}
}

    \item En un estudio que relaciona el promedio de calificaciones en la universidad con el
tiempo dedicado a diversas actividades, usted reparte un cuestionario a varios estudiantes
en el que les pregunta cuántas horas destinan semanalmente a estas cuatro categorías:
estudiar, dormir, trabajar y tiempo libre. Cualquier actividad se sitúa en una de las
categorías, de modo que la suma de horas de cada estudiante en las cuatro debe ser 168.

\begin{enumerate}
    \item[(i)] (5 puntos) En el modelo, explique por qué el modelo infringe la suposición de no Colinealidad Perfecta.

    \[
        GPA = \beta_0 + \beta_1 \, estud + \beta_2 \, dorm + \beta_3 \, trab + \beta_4 \, tlib + u,
    \]
\color{blue}
La suposición (no colinealidad perfecta) exige que no exista una relación lineal exacta entre los regresores. Sin embargo, en este caso se verifica la restricción
    \[
        estud + dorm + trab + tlib = 168,
    \]
    que puede reescribirse, por ejemplo, como
    \[
        tlib = 168 - estud - dorm - trab.
    \]
    Por lo tanto, $tlib$ es una combinación lineal exacta de las otras tres variables y de una constante. Esto implica que los regresores son linealmente dependientes y la matriz de diseño no tiene rango completo. En consecuencia, el modelo infringe la suposición (no colinealidad perfecta).
\color{black}

    \item[(ii)] (5 puntos) ¿Cómo reformularía el modelo para que sus parámetros tuvieran una interpretación  útil y satisficieran la suposición de Colinealidad Imperfecta?
\color{blue}
Para que los parámetros tengan una interpretación útil y se satisfaga no colinealidad perfecta, se puede eliminar una de las cuatro variables. 
    \[
        GPA = \beta_0 + \beta_1 estud + \beta_2 dorm + \beta_3 trab + \beta_4 tlib + u
    \]
    y sustituimos:
    \[
        GPA = \beta_0 + \beta_1 estud + \beta_2 dorm + \beta_3 trab + \beta_4 (168 - estud - dorm - trab) + u.
    \]
    Reagrupando términos:
    \[
        GPA = (\beta_0 + 168\beta_4)
        + (\beta_1 - \beta_4) estud
        + (\beta_2 - \beta_4) dorm
        + (\beta_3 - \beta_4) trab
        + u.
    \]
    Definiendo nuevos parámetros:
    \[
        \alpha_0 = \beta_0 + 168\beta_4,\quad
        \alpha_1 = \beta_1 - \beta_4,\quad
        \alpha_2 = \beta_2 - \beta_4,\quad
        \alpha_3 = \beta_3 - \beta_4,
    \]
    el modelo puede escribirse como
    \[
        GPA = \alpha_0 + \alpha_1 estud + \alpha_2 dorm + \alpha_3 trab + u,
    \]
    sin incluir $tlib$ como regresor. En este modelo ya no hay colinealidad perfecta entre las variables explicativas, por lo que se satisface no colinealidad perfecta. Además, cada coeficiente $\alpha_j$ se interpreta como el efecto de aumentar en una hora la actividad correspondiente, reduciendo en una hora el tiempo libre (que actúa como categoría omitida), manteniendo constantes las otras dos actividades.
\color{black}
\end{enumerate}

    
    \item Si en un modelo existen problemas por variable omitida y este se puede solucionar con una variable instrumental, no necesito saber cuál es la variable omitida, sólo que existe una. Comente.\\
    {\color{blue}
    \textbf{Respuesta:} Falso. Siempre es necesario saber cuál es la variable omitida, dado que hay que testear que el instrumento (Z) sea exógeno, es decir, que no este correlacionado con el término de error. Si hay una variable omitida, esta está en el término de error, y el instrumento tiene que no estar correlacionada con dicha variable, por lo cual es necesario saber a qué corresponde esa variable (supuesto de Exogeneidad). Además, necesito que el instrumento cumpla con el supuesto de Relevancia, es decir que la variable Z correlacione con la variable omitida X, por lo cual, nuevamente, es necesario saber a qué corresponde esta última.}
    
    \item Un investigador ha simulado tres series temporales de 200 observaciones. 
Cada una corresponde a uno de los siguientes procesos:

\begin{itemize}
    \item Un proceso de \textbf{ruido blanco} (white noise).
    \item Un proceso \textbf{MA(1)} con parámetro $\theta = 0.8$:
    \[
        X_t = \varepsilon_t + 0.8\,\varepsilon_{t-1},
    \]
    \item Un proceso \textbf{AR(1)} con parámetro $\phi = 0.8$:
    \[
        X_t = 0.8\,X_{t-1} + \varepsilon_t,
    \]
\end{itemize}

donde en todos los casos $\varepsilon_t \sim \text{i.i.d. } N(0,1)$.

Para cada caso se presentan dos gráficos:
\begin{itemize}
    \item A la izquierda: la serie temporal simulada $\{X_t\}$.
    \item A la derecha: su función de autocorrelación (ACF) hasta el rezago 30, con bandas de confianza aproximadas.
\end{itemize}

Los gráficos han sido rotulados como \textbf{Caso A}, \textbf{Caso B} y \textbf{Caso C} (sin indicar el tipo de proceso que representan).

\begin{figure}[H]
    \centering
    \includegraphics[width=0.65\linewidth]{certamenes//certamenesDiurno25-2/seriesACF.png}
\end{figure}



\begin{enumerate}
    \item[(a)] (4 puntos) Identifique cuál de los tres casos corresponde al proceso de \textbf{ruido blanco}. Justifique su respuesta con base en el comportamiento de la serie y en la forma de su ACF.

{\color{blue}
    El ruido blanco es el Caso A. (2 puntos) Para todos los regazos desde 1 en adelante la función de autocoreelación tiene el valor de 0, característica de este tipo de proceso estocástico que es iid. (justificacion 2 puntos) }

    \item[(b)] (6 puntos) Para los otros dos casos, indique cuál corresponde al proceso \textbf{MA(1)} con $\theta = 0.8$ y cuál corresponde al proceso \textbf{AR(1)} con $\phi = 0.8$. Justifique su elección explicando cómo se refleja en la ACF el hecho de que:
    \begin{itemize}
        \item En el MA(1) la autocorrelación se “corta” después de un número finito de rezagos.
        \item En el AR(1) la autocorrelación decae gradualmente de forma geométrica.
    \end{itemize}

 {\color{blue}   El proceso MA(1) es el caso B (1 punto), vemos que su función de autocorrelación se corta en el primer rezago tal como se señala en el enunciado, su función de autocorrelacion tendra valores distintos de 0 solo para el primer rezago.

    Para el proceso MA(1)
\[
X_t = \varepsilon_t + \theta \varepsilon_{t-1}, 
\qquad \varepsilon_t \sim \text{i.i.d. }(0,\sigma^2),
\]
la autocorrelación de primer orden se obtiene de la siguiente manera.

\textbf{Varianza}:
\[
\gamma(0) = \operatorname{Var}(X_t)
= \sigma^2(1+\theta^2).
\]

\textbf{Covarianza con el primer rezago}:
\[
\gamma(1) 
= \operatorname{Cov}(X_t, X_{t-1})
= \theta \sigma^2.
\]

\textbf{Autocorrelación de primer orden}:
\[
\rho(1) 
= \frac{\gamma(1)}{\gamma(0)}
= \frac{\theta \sigma^2}{\sigma^2 (1+\theta^2)}
= \frac{\theta}{1+\theta^2}.
\]

Para \(\theta = 0.8\):
\[
\rho(1) = \frac{0.8}{1 + 0.8^2}
= \frac{0.8}{1.64}
\approx 0.488.
\]

No es necesario todos los calculos para el puntaje basta explicar. (3 puntos por el MA)
\vspace{0.5cm}

{\bf Para el proceso AR(1):}
\[
X_t = \phi X_{t-1} + \varepsilon_t,
\qquad \varepsilon_t \sim \text{i.i.d. } (0,\sigma^2), 
\qquad |\phi| < 1.
\]


\[
\gamma(0) = \operatorname{Var}(X_t).
\]
Usando la ecuación del modelo:
\[
X_t = \phi X_{t-1} + \varepsilon_t,
\]
obtenemos
\[
\gamma(0) 
= \operatorname{Var}(\phi X_{t-1} + \varepsilon_t)
= \phi^2 \operatorname{Var}(X_{t-1}) + \operatorname{Var}(\varepsilon_t)
= \phi^2 \gamma(0) + \sigma^2.
\]
Despejando:
\[
\gamma(0)(1 - \phi^2) = \sigma^2
\quad\Rightarrow\quad
\gamma(0) = \frac{\sigma^2}{1-\phi^2}.
\]

Autocovarianza:

\[
\gamma(1) = \operatorname{Cov}(X_t, X_{t-1}).
\]
Usando de nuevo el modelo:
\[
\gamma(1) 
= \operatorname{Cov}(\phi X_{t-1} + \varepsilon_t,\, X_{t-1})
= \phi\,\operatorname{Cov}(X_{t-1}, X_{t-1})
    + \operatorname{Cov}(\varepsilon_t, X_{t-1}).
\]
Como $\varepsilon_t$ es ruido blanco e independiente de $X_{t-1}$:
\[
\operatorname{Cov}(\varepsilon_t, X_{t-1}) = 0,
\]
por lo que
\[
\gamma(1) = \phi\,\gamma(0).
\]

Entonces la \textbf{autocorrelación de primer orden} es:
\[
\rho(1) = \frac{\gamma(1)}{\gamma(0)} = \frac{\phi\,\gamma(0)}{\gamma(0)} = \phi.
\]



De manera análoga, para $k \geq 1$:
\[
\gamma(k) = \operatorname{Cov}(X_t, X_{t-k})
= \operatorname{Cov}(\phi X_{t-1} + \varepsilon_t,\, X_{t-k})
= \phi\,\operatorname{Cov}(X_{t-1}, X_{t-k})
= \phi\,\gamma(k-1).
\]

Esta relación recursiva implica
\[
\gamma(k) = \phi^k \gamma(0), \qquad k = 0,1,2,\dots
\]
y, por tanto, la \textbf{función de autocorrelación} es
\[
\rho(k) = \frac{\gamma(k)}{\gamma(0)} = \frac{\phi^k \gamma(0)}{\gamma(0)} = \phi^k,
\qquad k = 0,1,2,\dots
\]

Es decir, en un AR(1) estacionario la autocorrelación de orden $k$ decae geométricamente como $\phi^k$.

No es necesario todo el calculo solo explicar porque la correlación con el  primer rezago es 0.8 y luego decrece. (3 puntos)
}
\end{enumerate}

    \item (10 puntos) En la presentación de Ph.D. Cristóbal Otero de su investigación \textit{Physicians’ Occupational Licensing and the Quantity-Quality trade-off} se distinguen dos efectos causales, en acceso a la salud y en mortalidad. ¿Qué estrategia empírica utilizan los autores para determinar la causalidad? ¿Cómo son esos efectos en términos de significancia y dirección?\\

    {\color{blue} 
    Los autores utilizan una variación quasi-experimental para aislar efecto de la cantidad y calidad de médicos en resultados de hospitales públicos. Específicamente utilizan Shift-Share IV, ya que:
    \begin{itemize}
        \item Migración y expansión de carreras de medicina en Chile generaron cambios exógenos en cantidad y calidad de médicos.
        \item Impacto distinto según la exposición de cada hospital
    \end{itemize}
    
    Esto permite estimar efecto causal de cantidad y calidad de médicos sobre resultados, debido a que un cambio en número de médicos de un hospital está impulsado por la migración externa y patrones de graduación, y no por calidad propia del hospital. \textbf{(5 puntos)}\\

    En términos de resultados \textbf{(1,25 puntos cada una)}:
    \begin{itemize}
        \item Acceso:
        \begin{itemize}
            \item Cantidad de médicos tiene un impacto causal y signifiativo en el acceso, positivo.
            \item Calidad de médicos no tiene un impacto significativo en acceso
        \end{itemize}
        
        \item Mortalidad:
        \begin{itemize}
            \item Cantidad de médicos tiene un impacto causal y signifiativo en el acceso, negativo.
            \item Calidad de médicos tiene un impacto causal y signifiativo en el acceso, negativo.
        \end{itemize}
    \end{itemize}
    }
\end{enumerate}

\section{RdeA9.1 Ejercicios (60 puntos)}

\begin{enumerate}
    \item (20 puntos) 
Considere el siguiente modelo que busca explicar el impacto innovador de un proyecto tecnológico,
medido mediante el índice \textit{ImpactScore}:

\[
ImpactScore_i
= \beta_0 
+ \beta_1\, Invest_i
+ \beta_2\, TeamExp_i
+ \beta_3\, Collab_i
+ u_i,
\]

donde:
\begin{itemize}
    \item $Invest_i$: inversión inicial destinada al proyecto (en miles de dólares),
    \item $TeamExp_i$: años totales de experiencia acumulada por el equipo fundador,
    \item $Collab_i$: número de alianzas con universidades, laboratorios u otras empresas.
\end{itemize}

La organización desea investigar si la relación entre estas variables es igual para dos tipos de startups:
\textit{startups verdes} (grupo 1) y \textit{startups de inteligencia artificial} (grupo 2). Para ello realiza un Test de Chow.

Los resultados obtenidos son:
\[
SSR_1 = 210, \qquad SSR_2 = 250, \qquad SSR_r = 500,
\]
con $n = 120$ observaciones totales.

Responda:

\begin{enumerate}
    \item[(a)] (5 puntos) Formule las hipótesis nula y alternativa del Test de Chow.


 \color{blue}
Las hipótesis del Test de Chow son:
    
    \[
    H_0:\ 
    (\beta_{0,1},\beta_{1,1},\beta_{2,1},\beta_{3,1})
    =
    (\beta_{0,2},\beta_{1,2},\beta_{2,2},\beta_{3,2}),
    \]
    es decir, los parámetros del modelo son iguales para startups verdes y startups de inteligencia artificial.
    
    \[
    H_1:\ 
    (\beta_{0,1},\beta_{1,1},\beta_{2,1},\beta_{3,1})
    \neq
    (\beta_{0,2},\beta_{1,2},\beta_{2,2},\beta_{3,2}),
    \]
    al menos uno de los parámetros difiere entre los dos grupos (existe un quiebre estructural).

\color{black}

    \item[(b)] (10 puntos) Indique si debe rechazarse la hipótesis nula. (considere el valor crítico al 5\%,
    $F_{0.05}(4,112)=2.45$.
\color{blue}

\[
    SSR_{ur} = SSR_1 + SSR_2 = 210 + 250 = 460.
    \]
    
    Utilizamos la fórmula del estadístico $F$ de Chow solicitada:
    
    \[
    F=
    \frac{(SSR_r - SSR_{ur}) / \big[(G-1)(k+1)\big]}
         {SSR_{ur} / \big[n - G(k+1)\big]},
    \]
    
    donde $SSR_r = 500$, $G = 2$, $n = 120$ y $k+1 = 4$ parámetros (incluyendo intercepto).
    
    \[
    (SSR_r - SSR_{ur}) = 500 - 460 = 40,
    \]
    \[
    (G-1)(k+1) = (2-1)(4) = 4,
    \]
    \[
    \frac{SSR_r - SSR_{ur}}{(G-1)(k+1)} = \frac{40}{4} = 10.
    \]
    
    Para el denominador:
    \[
    n - G(k+1) = 120 - 2(4) = 120 - 8 = 112,
    \]
    \[
    \frac{SSR_{ur}}{n - G(k+1)} = \frac{460}{112} \approx 4.107.
    \]
    
    Entonces:
    \[
    F \approx \frac{10}{4.107} \approx 2.43.
    \]

    Dado que:
    \[
    F_{\text{calc}} = 2.43 < F_{0.05}(4,112)=2.45,
    \]
    no se rechaza la hipótesis nula al 5\%.
\color{black}
    
    \item[(c)] (5 puntos) Interprete el resultado en el contexto de la innovación tecnológica.
\end{enumerate}
\color{blue}
En el contexto de la innovación tecnológica, al no rechazar $H_0$ concluimos que
    no existe evidencia estadística de que los coeficientes del modelo sean distintos entre
    startups verdes y startups de inteligencia artificial. Por lo tanto, la relación entre
    \textit{ImpactScore}, la inversión, la experiencia del equipo y las colaboraciones
    no presenta diferencias significativas entre ambos tipos de startups.
    \color{black}
    \item (20 puntos) Se dispone de la siguiente serie del PIB real de Chile, medida en millones de pesos a precios constantes, para el período 2000--2024:

\begin{center}
    \begin{figure}[H]
        \centering
        \includegraphics[width=0.7\linewidth]{certamenes/Certamenes T2/imagen pib.png}
        \caption{PIB real de Chile (2000--2024)}
        \label{fig:pib_chile}
    \end{figure}
\end{center}

Un investigador está interesado en analizar las propiedades dinámicas del PIB real. 
\begin{enumerate}
    \item[(a)]  (8 puntos) Suponiendo que el investigador plantea que el PIB real podría ser una serie integrada de orden 1, escriba el modelo de regresión auxiliar del test Dickey--Fuller  correspondiente. Indique claramente la forma funcional, la variable dependiente y el término cuyo coeficiente es relevante para el contraste de raíz unitaria.


{\color{blue}
El modelo apropiado en este caso es el del test Dickey--Fuller que incluye una tendencia 
lineal:

\begin{eqnarray}
 y_t &=& \alpha + \rho y_{t-1} + \beta t + \varepsilon_t\\
 y_t - y_{t-1} &=& \alpha + \rho y_{t-1} - y_{t-1} + \beta t + \varepsilon_t\\
 \Delta y_t &=& \alpha + (\rho - 1) y_{t-1} + \beta t + \varepsilon_t\\
 \Delta y_t &=& \alpha + \phi y_{t-1} + \beta t + \varepsilon_t,
\end{eqnarray}

donde $\Delta y_t = y_t - y_{t-1}$ es la primera diferencia del PIB real.

La \textbf{regresión auxiliar} del test Dickey--Fuller con tendencia lineal queda entonces dada por:
\[
\Delta y_t = \alpha + \phi y_{t-1} + \beta t + \varepsilon_t .
\]

En esta regresión:
\begin{itemize}
    \item La \textbf{variable dependiente} es $\Delta y_t$, la primera diferencia del PIB real.
    \item Los \textbf{regresores} son el rezago $y_{t-1}$ y la tendencia determinística $t$ (además del intercepto $\alpha$).
    \item El \textbf{coeficiente relevante para el contraste de raíz unitaria} es $\phi$, asociado al término $y_{t-1}$.
\end{itemize}
}
       
    \item[(b)] (4 puntos) Formule la hipótesis nula y la hipótesis alternativa del test Dickey--Fuller, utilizado para evaluar la presencia de una raíz uonitaria en el PIB real.

{\color{blue}
     \begin{itemize}
            \item $H_0: \phi = 0 \implies \rho = 1 \implies \text{Raíz unitaria }$
            \item $H_1: \phi < 0 \implies \rho < 1 \implies \text{Serie estacionaria }$
        \end{itemize}}

    \item[(c)] (8 puntos) A continuación se presentan los resultados de la estimación del modelo ADF descrito en el inciso (a):

    \begin{center}
        \begin{figure}[H]
            \centering
            \includegraphics[width=0.8\linewidth]{certamenes/Certamenes T2/rlst.png}
            \caption{Resultados del test ADF para el PIB real}
            \label{fig:adf_pib}
        \end{figure}
    \end{center}

   donde L(pib,1) representa el rezago de orden 1 del PIB. Utilizando estos resultados y los valores críticos apropiados del test Dickey--Fuller al 5\% de significancia, determine si existe raíz unitaria en el PIB real de Chile. 


{\color{blue}  El valor del estadístico t para el rezago del PIB real es:

$$t_{DF}={-0.5208\over 0.1777}=-2.931$$ 

Considerando que se cuenta con 25 observaciones, el valor crítico de la tabla al 5\% y con tendencia es -3.6. Luego, como 2.931<3.6, no se rechaza la hipótesis nula. El PIB real presenta una raíz unitaria.   }

   
\end{enumerate}




    \item (20 puntos)  Usando datos sobre gasto en transporte se estimó la siguiente ecuación para trabajadores de la región metropolitana


{\scriptsize
\begin{equation*}
\begin{array}{ll}
\widehat{GT} = & 1,2 \quad + \quad0,1 \cdot dist\quad + \quad 0,04 \cdot auto \quad - \quad 0,03 \cdot tele
 \\
& (0,5)  \hspace{0,8cm} (0,05) \hspace{1,3cm} (0,019) \hspace{1,3cm} (0,006)\\ 
& [0,5]  \hspace{0,9cm} [0,0511] \hspace{1,2cm} [0,02] \hspace{1,4cm}  [0,004] \\
& n = \text{450} \quad R^2 = \text{0,73}
\end{array}
\end{equation*}}



Aquí $GT$ es el gasto promedio en transporte en un mes que realiza un individuo; $dist$ es la distancia al trabajo; $auto$ cantidad de días que el individuo utiliza automóvil para ir al trabajo en el mes; $tele$ cuántos días de teletrabajo tiene el individuo en un mes. Entre paréntesis y entre corchetes se dan, respectivamente,
los errores estándar usuales y los errores estándar robustos a la heterocedasticidad.


\begin{enumerate}
    \item \textbf{(6 puntos)} ¿Tienen las variables  $dist$, $auto$ y $tele$ los efectos estimados esperados? ¿Cuáles de estas variables son estadísticamente significativas al nivel de 5\% suponiendo que se cumple el supuesto de homocedasticidad?

{\color{blue}
\textbf{(1 punto cada una)}:
\begin{itemize}
    \item \(dist\): El coeficiente es \(0,1\), positivo, lo cual sugiere que un mejor promedio acumulado previo está asociado con un mayor \(GT\). Esto es coherente ya que una mayor distancia al trabajo debería asociarse con un mayor gasto en transporte.
    \item \(auto\): El coeficiente es \(0,04\), positivo pero de pequeña magnitud, indicando que el uso de automóvil incrementa el gasto promedio de transporte.
    \item \(tele\): El coeficiente es \(-0,03\), negativo pero de pequeña magnitud, indicando que un aumento en un día de teletrabajo disminuye el gasto en transporte (en 0,03 unidades monetarias).
\end{itemize}



Para evaluar la significancia estadística de estas variables, calculamos los valores \(t\) como:

\[
t = \frac{\text{Coeficiente}}{\text{Error estándar}}
\]

Usando errores estándar usuales \textbf{(1 punto cada una)}:
\begin{itemize}
    \item \(dist\): \[
    t = \frac{0.1}{0.05} = 2 \quad (\text{Mayor que } 1.96 \rightarrow \text{Significativo})
    \]
    \item \(auto\): \[
    t = \frac{0.04}{0.019} = 2.11 \quad (\text{Mayor que } 1.96 \rightarrow \text{Significativo})
    \]
    \item \(tele\): \[
    t = \frac{0.03}{0.006} = 5 \quad (\text{Mayor que } 1.96 \rightarrow \text{Significativo})
    \]
\end{itemize}
}

\item \textbf{(8 puntos)} Suponga que se sospecha heterocedasticidad, ¿se sigue manteniendo la significancia estadística al 5\%? Concluya respecto a lo obtenido.

{\color{blue}
Usando errores estándar robustos \textbf{(1 punto cada una)}:
\begin{itemize}
    \item \(dist\): \[
    t = \frac{0.1}{0.0511} = 1.957 \quad (\text{Menor que } 1.96 \rightarrow \text{No significativo})
    \]
    \item \(auto\): \[
    t = \frac{0.04}{0.02} = 2 \quad (\text{Mayor que } 1.96 \rightarrow \text{Significativo})
    \]
    \item \(tele\): \[
    t = \frac{0.03}{0.004} = 7.5 \quad (\text{Mayor que } 1.96 \rightarrow \text{Significativo})
    \]
\end{itemize}

Utilizando errores estándar robustos la variable distancia al trabajo no es estadísticamente distinta de cero, puesto que su $t \; estadístico$ es menor que 1,96, mientras que $auto$ y $tele$ sí lo son. \textbf{(3 puntos)}. 
}

\item \textbf{(4 puntos)} Plantee el test de Breusch-Pagan para heterocedasticidad. Sabiendo que el $R_{aux}^2$ de la regresión auxiliar alcanza el valor de 0.1, ¿es posible afirmar que hay heterocedasticidad?

{\color{blue} El test de Breusch-Pagan viene de correr la regresión:

$$\hat{u}_i = \delta_0 + \delta_1 \cdot dist + \delta_2 \cdot auto + \delta_3 \cdot tele$$

Luego \textbf{(2 puntos)}, 

$$H_0: \delta_1 = \delta_2 = \delta_3 = 0$$
$$H_0: \text{al menos un } \delta_j \neq 0$$

Con la información del enunciado,e puede calcular el test LM que es equivalente a Breusch-Pagan de la forma \textbf{(1 punto)}:

$$LM: n R_{aux}^2 \sim \chi^2_{(k)}$$
$$LM: 450 \cdot 0.1 \sim 7.81$$

Como $45 > 7.81$, entonces se rechaza la hipótesis nula y se concluye que hay heterocedasticidad \textbf{(1 punto)}.}

\item \textbf{(4 puntos)} Se sospecha que puede haber multicolinealidad entre las variables, por lo que se corre la siguiente regresión:

$$auto = \delta_0 + \delta_1 \cdot dist - \delta_2 \cdot tele$$

Obteniendo un $R_j^2 = 0.89$. Con dicha información, ¿se puede afirmar que hay multicolinealidad?

{\color{blue}

No, no es posible afirmarlo, puedo que si calculamos el VIF:

$$VIF = \frac{1}{1 - R_j^2}$$

Siendo $R_j^2$ el coeficiente de determinación de la regresión auxiliar. Reemplazando:

$$VIF = \frac{1}{1 - 0.89} = \frac{1}{0.1} \approx 9.09 < 10$$

$\implies$ No se puede afirmar que exista multicolinealidad, aunque el $R_a^2$ sea elevado.}

\end{enumerate}
\end{enumerate}
\end{document}